% Source: https://artofproblemsolving.com/community/c3h2834449_geometry_amc_problem
% Stable identity: aops-topic-2834449-q01
Since $YE\parallel AD$, triangles $XYE$ and $XDA$ are similar, so $\dfrac{XE}{XA}=\dfrac{XY}{XD}$. Since $EF\parallel AC$, triangles $XEF$ and $XAC$ are similar, so $\dfrac{XF}{XC}=\dfrac{XE}{XA}=\dfrac{XY}{XD}$. By the power of the point $X$, $XC\cdot XG=XB\cdot XD$. Hence
\[XF\cdot XG=\dfrac{XY}{XD}\cdot XC\cdot XG=XY\cdot XB.\]
\begin{asy}
usepackage("amsmath");
size(4.5cm);
defaultpen(fontsize(10pt));
pair A=dir(200),B=dir(120),C=dir(75),D=dir(-30);
draw(circle((0,0),1),gray);draw(A--B--C--D--cycle);draw(B--D);draw(A--C,dotted);
pair X=D+(B-D)/4,Y=B+(D-B)*11/36;dot(X);dot(Y);
pair G=intersectionpoints(C--(C+3*(X-C)),circle((0,0),1))[0];draw(C--G,dashed);dot(G);
label("$A$",A,W);label("$B$",B,NW);label("$C$",C,NE);label("$D$",D,SE);label("$X$",X,NE);label("$Y$",Y,NE);label("$G$",G,S);
\end{asy}

With $BD=t$: $XB=\dfrac34t$ and $XY=t-\dfrac t4-\dfrac{11t}{36}=\dfrac49t$, so $XF\cdot XG=\dfrac{t^2}3$.

Opposite angles $A$ and $C$ are supplementary, so by the law of cosines $t^2=9+64-48\cos A=4+36+24\cos A$. Then $72\cos A=33$, $\cos A=\dfrac{11}{24}$, and $t^2=73-22=51$. So $XF\cdot XG=\boxed{\textbf{(A)}\ 17}$.
